% =============================================================================
% 8. THEORETICAL COMPARISON WITH RELATED APPROACHES
% Target: 1 page (~500 words)
% Source: cap4 sec:s4-comparison (tab:mc-comparison) + extended discussion
% Status: SKELETON
%
% Purpose: position PATHCAST against three classes of related approaches
% on EXPRESSIVENESS and ASSUMPTIONS, not on empirical numbers.
% =============================================================================

\section{Theoretical Comparison with Related Approaches}
\label{sec:comparison}

\textbf{[TODO — 1 p.]}
We compare PATHCAST with three families of related approaches along
five theoretical dimensions: process model, sampling unit, rework
modeling, state-specific delays, and distributional outputs.

\begin{table}[ht]
\centering
\caption{Theoretical comparison of forecasting approaches along five
dimensions. ``--'' indicates the approach does not capture the
dimension.}
\label{tab:approach-comparison}
\small
\begin{tabular}{@{}l l l l l@{}}
\toprule
\textbf{Dimension}    &
\textbf{Throughput MC}~\cite{vacanti2015} &
\textbf{Markov analyt.}~\cite{kemeny1976finite} &
\textbf{NM-SPN}~\cite{RoggeSolti2015NonMarkovianSPN} &
\textbf{PATHCAST} \\
\midrule
Process model         & --                 & exogenous          & Petri net          & PM-discovered \\
Sampling unit         & daily throughput   & none (analytical)  & timed firing       & transition + sojourn \\
Rework                & --                 & via $N$           & via Petri loops    & via $P_{ab}$ + $\hat{F}$ \\
State-specific delay  & --                 & via $\bar{d}_s$    & general dist.      & empirical $\hat{F}_s$ \\
Distributional output & completion only    & --                 & lead-time only     & LT + outcome + path + WIP \\
\bottomrule
\end{tabular}
\end{table}

\paragraph{Throughput Monte Carlo.}
\textbf{[TODO — 2 sentences.]}
Throughput MC samples future cumulative completion from the empirical
throughput distribution of a team. It is operationally simple but
process-blind: rework, state-specific delays, and path information are
not represented. By Theorem~\ref{thm:information}, throughput MC
cannot exploit the information gain $I(T;\psi)$ encoded in the
trajectory.

\paragraph{Analytical absorbing Markov chains.}
\textbf{[TODO — 2 sentences.]}
The closed-form expectations $N\mathbf{1}$ and $NR$ are
deterministic; they yield expected lead time but not a forecast
distribution. PATHCAST uses these expressions for analytical
cross-validation (Section~\ref{sec:numerical-example}) but produces
distributional forecasts via Monte Carlo, exposing percentiles and tail
risk that the analytical baseline cannot.

\paragraph{Non-Markovian stochastic Petri nets (NM-SPN).}
\textbf{[TODO — 2 sentences.]}
NM-SPN~\cite{RoggeSolti2015NonMarkovianSPN} parameterize timed firings
with general distributions over an exogenously specified Petri net.
PATHCAST differs in two respects: (i) the process model is
\emph{discovered} from the event log rather than supplied externally,
removing a strong modelling burden; (ii) sojourn distributions are
empirical rather than parametric, avoiding misspecification of the
heavy-tailed, zero-inflated distributions typical of software
processes.

\paragraph{Process-prediction approaches.}
\textbf{[TODO — 2 sentences.]}
The predictive process monitoring
literature~\cite{polato2018time, verenich2019survey} addresses remaining
time prediction in business processes, typically via supervised
learning. PATHCAST occupies a complementary niche: it integrates a
stochastic process model (not a black-box predictor) and is specialised
to the heterogeneous data sources of the SDLC. The ML component
(Section~\ref{sec:ml-component}) bridges to this literature by adding
case-level point forecasts on top of the population-level distribution.
